% figures/ch04-dirichlet-character.tex
% Dirichlet 特征 χ mod 5 的值示意表
\begin{figure}[ht]
\centering
\begin{tikzpicture}[font=\small, >=Stealth]

% === 标题 ===
\node[font=\normalsize\bfseries] at (0, 2.8) {$(\Z/5\Z)^\times$ 上的特征};

% 单位根在单位圆上的位置
\draw[gray!40] (0,0) circle (1.5cm);
\draw[->] (-2, 0) -- (2, 0) node[right, font=\footnotesize] {$\Re$};
\draw[->] (0, -2) -- (0, 2) node[above, font=\footnotesize] {$\Im$};

% 四个四次单位根
\foreach \ang/\val/\lab in {
  0/{$1$}/{$\chi(1)=1$},
  90/{$i$}/{$\chi(2)=i$},
  180/{$-1$}/{$\chi(3)=-1$},
  270/{$-i$}/{$\chi(4)=-i$}
}{
  \fill[blue!70!black] ({1.5*cos(\ang)}, {1.5*sin(\ang)}) circle (3pt);
  \node[font=\footnotesize] at ({1.9*cos(\ang)}, {1.9*sin(\ang)}) {\val};
}

% 对应的 a mod 5
\node[red!70!black, font=\footnotesize] at (1.5+0.4, -0.3) {$a=1$};
\node[red!70!black, font=\footnotesize] at (0.2, 1.5+0.3) {$a=2$};
\node[red!70!black, font=\footnotesize] at (-1.5-0.5, -0.3) {$a=3$};
\node[red!70!black, font=\footnotesize] at (0.2, -1.5-0.3) {$a=4$};

% 生成元 2 的标注
\draw[->, thick, orange!70!black, bend left=20]
  ({1.5*cos(0)+0.1}, {1.5*sin(0)}) to ({1.5*cos(90)}, {1.5*sin(90)+0.1});
\node[orange!70!black, font=\footnotesize] at (1.2, 0.9) {$\times 2$};

% === 右侧：表格 ===
\begin{scope}[xshift=5cm, yshift=0.3cm]
\node[font=\footnotesize\bfseries] at (0, 2.2) {特征值表};

% 表格内容
\foreach \row/\a/\ca/\cb/\cc/\cd in {
  0/{$a \bmod 5$}/{$1$}/{$2$}/{$3$}/{$4$},
  -0.7/{$\chi_0$ (主特征)}/{$1$}/{$1$}/{$1$}/{$1$},
  -1.4/{$\chi_1$}/{$1$}/{$i$}/{$-i$}/{$-1$},
  -2.1/{$\chi_2$}/{$1$}/{$-1$}/{$-1$}/{$1$},
  -2.8/{$\bar\chi_1$}/{$1$}/{$-i$}/{$i$}/{$-1$}
}{
  \node[font=\footnotesize, anchor=east] at (-1.2, \row) {\ca};
  \node[font=\footnotesize] at (-0.4, \row) {\cb};
  \node[font=\footnotesize] at (0.4, \row) {\cc};
  \node[font=\footnotesize] at (1.2, \row) {\cd};
  \node[font=\footnotesize, anchor=east] at (-2.0, \row) {\a};
}
% 分隔线
\draw[gray] (-2.3, -0.35) -- (1.5, -0.35);
\draw[gray] (-2.3, 1.8) -- (1.5, 1.8);

% 列标签
\node[font=\footnotesize\bfseries, anchor=east] at (-2.0, 1.9) {特征 $\backslash$ $a$};

\end{scope}

\end{tikzpicture}
\caption{模 $5$ 的 Dirichlet 特征。$(\Z/5\Z)^\times = \{1,2,3,4\}$
是循环群，生成元 $2$（因为 $2^1=2, 2^2=4, 2^3=3, 2^4=1 \pmod 5$）。
共有 $4$ 个特征，对应四次单位根 $\{1, i, -1, -i\}$（左侧单位圆）。
$\chi_0$ 是主特征，$\chi_1$ 是将生成元 $2$ 映到 $i$ 的特征。
所有特征延拓到 $\Z$ 上时，令 $\chi(n) = 0$ 若 $\gcd(n,5)>1$。}
\label{fig:dirichlet-character}
\end{figure}
